\chapter*{Declarations}
\markboth{Declarations}{Declarations}
\addcontentsline{toc}{chapter}{Declarations}

\section*{Use of AI Tools}

I acknowledge the use of Claude (Anthropic, \url{https://claude.ai/}) and GPT models from OpenAI, including \texttt{GPT-4.1-mini}, \texttt{GPT-4o-mini}, \texttt{GPT-4o}, and \texttt{GPT-4.5} (OpenAI, \url{https://platform.openai.com/}), during this project. Claude was used to provide feedback on portions of this report, to assist with debugging Python code, and to clarify understanding of papers read during the course of the project. Several diagrams in this report were initially hand-drawn by me and subsequently refined into clean vector-style figures using ChatGPT (\texttt{GPT-4.5}) via the web interface.\par

GPT models were used to generate the emotion-rewrite dataset described in Chapter~\ref{chap:development1_emo_steering}, producing affective variants of source utterances across Plutchik's eight basic emotions at three intensity levels. They were also used to evaluate the outputs of the steering experiments, rating the emotional intensity and texture shifts produced by interventions along the extracted vectors. This use is described explicitly in the methodology (Sections~\ref{sec:steering_experiments},~\ref{sec:ch2-causal-validation}). No content generated by AI has been presented as my own work. \par

\section*{Ethical Considerations}

This project does not involve human subjects in the traditional sense, but it does include a human evaluation study in which participants rated generated text for affective qualities. Participation was voluntary, no personally identifiable information was collected, and the study was conducted with the informed consent of all participants.\par

More broadly, the project contributes to the goal of making affective manipulation in language models interpretable and controllable. Systems that steer the emotional tone of generated text raise questions about potential misuse, for instance in generating manipulative or deceptive content at scale. The methods developed here operate at the level of interpretable geometric components, a property that supports auditability and transparency rather than obfuscation. Nonetheless, the dual-use potential of affective steering technology warrants continued attention as such systems mature.\par

\section*{Sustainability}

The computational footprint of this project was modest. No training was performed, i.e.\ all experiments used frozen pretrained model weights. Activation extraction and steering experiments ran on a single machine, using CPU inference for the smaller models and a consumer GPU for \texttt{Llama-3.1-8B-Instruct}. The dataset construction step used \texttt{GPT-4.1-mini} via API to generate a large number of short rewrites; although no model training was performed, this remote computation constitutes the main external computational cost of the project.

\section*{Data and Materials}

All code for activation extraction, CAA vector construction, decomposition, pooled PCA, steering experiments, and human evaluation analysis is maintained in the canonical repository for this project:

\begin{itemize}
    \item \url{https://github.com/maplesugano/EmotionEngine_v2} --- the canonical repository for this report, containing all final experiments, analysis scripts, and the reproducibility checklist (Appendix~\ref{app:reproducibility})
\end{itemize}

Two earlier repositories exist for historical reference only and are not needed to reproduce any result reported here:
\begin{itemize}
    \item \url{https://github.com/maplesugano/EmotionEngine} --- an early-stage prototype preceding the current architecture
    \item \url{https://github.com/maplesugano/Dissertation} --- initial notes and exploratory code from the first weeks of the project
\end{itemize}